\chapter{Delivery and validation}\label{ch:delivering}

A compiled program is packaged into an image that Metal loads, and its result is accepted only after a ladder of
checks. Model graphs are then assembled from these programs.

\section{Image and validation ladder}\label{ch:image-result-validated}

\code{scanlink.author} writes every byte of the image: the Mach-O object, the metadata, the library and the binary
archive (MM~\mm{25.147}, MM~\mm{25.148}). The metadata holds the register count, the system registers, the
bindings and the tensor-unit enable. The census of every byte of 421 delivered archives, the 148 constants Apple's
loader requires and what Metal checks are in \cref{sec:object-library-archive}.

Validation is a ladder with the cheapest check first (\cref{tab:validation-ladder}; MM~\mm{23.6}).

\begin{xltabular}{\linewidth}{@{}l>{\hsize=1.34\hsize}X>{\hsize=0.66\hsize}X@{}}
\caption{The validation ladder in the order its checks run, with the tool that performs each.}\label{tab:validation-ladder}\\
\toprule
Step & Check & Where \\
\midrule
\endfirsthead
\tablecontinued{3}
\toprule
Step & Check & Where \\
\midrule
\endhead
1 & every instruction decodes back to what was asked for & the encoders \\*
2 & the whole body decodes, with one final \op{684}, no
R126+ operand and no read after release & \code{agxforge/g17/verify.py}, \code{releasecheck.py} \\*
3 & the contract and metadata agree with the code & \code{scanlink.author} \\*
4 & the CPU emulator runs the program from its bytes against a reference,
with the race check & \code{tools/g17emu.py} \\*
5 & one hardware dispatch per kernel, its output sentinel-filled, compared
with the reference & \code{tools/g17deliver.py} \\*
6 & Apple's compiler builds the same source, and the outputs
are compared with a control input that must differ & \code{tools/g17sourceexec.py} \\*
\bottomrule
\end{xltabular}

Steps 4 to 6 in more detail:

\begin{itemize}
\item Step 5 is exact, with one named exception. No CPU reference reproduces \op{1272}, so kernels that
  use it are checked against an enclosure with a wrong-base control. These are the decode and prefill attention in
  their \code{hw\_exp2} forms, and the fast SwiGLU (MM~\mm{25.144.5}, MM~\mm{25.183}). One input set is checked per kernel, not
  every input.
\item The emulator decodes the bytes, not the IR, so it reproduces the hardware from the program itself (MM~\mm{25.195},
  MM~\mm{25.197}).

  \begin{itemize}
  \item On shipped model graphs it is byte-identical to the GPU after every dispatch: one decode position of the
  batched graph (316 dispatches), one of single-stream decode (172), and the 1,024-token prefill through layers
  0 and 1.
  \item On Apple's own census programs it matches Apple's GPU output in 53 cases, with 0 differing and 45 refused by
  name.
  \item Its "wp" tier also admits semantics that only whole programs verified on hardware.
  \end{itemize}
\item Step 6 uses Apple's compiler only for comparison. Comparing outputs on the same source catches places where
  this compiler's reading of the language differs from Apple's. The cast in \cref{ch:selection-allocation-general} was found that way.
\end{itemize}

\section{Model graphs}\label{ch:model-graphs}

A model is compiled kernel by kernel and then assembled into a static graph (MM~\mm{25.138}, MM~\mm{25.141.15}):

\begin{itemize}
\item \code{tools/g17deliver.py build} compiles every kernel kind a model needs, authors each bundle, and dispatches it once.
  It writes the delivery index only if every check passes.
\item \code{tools/g17q4graph.py build} turns the model's weights and a config into \code{graph.json}, which holds:

  \begin{itemize}
  \item named arenas, one buffer each with guards, and their initial contents;
  \item the ordered dispatch list, with each binding's arena and offset;
  \item any per-token host writes.
  \end{itemize}
\item \code{tools/g17prefillgraph.py} adds the prefill section, which processes the prompt in M-row chunks before decoding
  starts (MM~\mm{25.142.10}).
\item A config selects each kernel's variant, so one model has many graphs (\code{tools/models/*.json}). They differ in their
  kernels, not in their structure.
\end{itemize}
